\documentclass[11pt]{article}
\usepackage[a4paper,margin=25mm]{geometry}
\usepackage{parskip}
\usepackage{hyperref}
\begin{document}
\begin{center}
{\Large\bfseries Machine Learning -- Contribution Information Sheet}\\[6pt]
{\bfseries CoRe-TFM: Selective Four-View Probability Reconciliation for Tabular Foundation Models}\\
Badampudi Agasthya Anirudh
\end{center}

\textbf{1. What is the main claim of the paper and why is it important?}

Recent work shows that tabular foundation models can produce direct marginal and conditional probabilities that cannot all arise from one joint distribution. The main claim is not that enforcing coherence is automatically beneficial. Rather, useful post-hoc correction requires \emph{selective four-view reconciliation}: trust in direct marginals versus conditional predictions must depend on their relative reliability, and validation should be allowed to retain an original factorization or simple pool when projection would impose a predictive cost. This matters for downstream uses that compose TFM probabilities into multivariate simulations, conditional queries, imputation, or uncertainty propagation.

\textbf{2. What precise evidence supports the claim?}

Two controlled protocols provide known conditional truth. Multiclass synthetic DGP sweeps vary dependence, sample size, feature dimension, class cardinality, and imbalance over 10 paired seeds while separating inner validation and outer testing. A new exact-view perturbation benchmark independently corrupts two direct marginals and two conditional families over six reliability regimes and 20 paired seeds. When direct marginals are reliable, Hard CoRe improves over arithmetic pooling by 0.0138 expected NLL on average (Holm-adjusted paired Wilcoxon $p<10^{-4}$); when marginals are noisy but conditionals are reliable, Hard CoRe is worse than arithmetic by 0.1423 NLL ($p<10^{-4}$). Directionally reliable regimes favor the corresponding original factorization, while symmetric noise often favors arithmetic/geometric pooling. These results support selective rather than unconditional reconciliation. The final real-TFM benchmark is pre-specified and remains a submission blocker until complete.

\textbf{3. What are the closest related contributions and how does this paper differ?}

The closest precursor is Kloetergens et al. (2026), \emph{Do Tabular Foundation Models Agree with Themselves?}, which defines marginalization/factorization consistency and documents violations across TFMs, but does not develop or evaluate a reconciliation method. DistPFN (Lee et al., 2026) performs post-hoc TFM probability adjustment for label shift, a different failure mode. Classical incompatible-conditional work and Bregman/Sinkhorn scaling literature provide general compatibility and optimization foundations; this manuscript does not claim those algorithms as new. The contribution is the TFM-specific four-view objective, hard/soft marginal-trust interpretation, selective deployment protocol, and controlled empirical characterization of the consistency--fidelity trade-off.

\textbf{4. Has any part of the contribution been published previously?}

No manuscript containing these CoRe-TFM methods or experiments has been published. The public GitHub repository is a research-code and working-draft repository. Any preprint status must be updated here if the manuscript is posted before journal submission.

\vfill
\textit{Working version: complete the real-TFM evidence and final author declarations before submission.}
\end{document}
